\documentclass{article}
\usepackage[margin=1in]{geometry}
\usepackage{tabularx}
\usepackage{array}
\usepackage{booktabs}
\usepackage{xcolor}
\usepackage{amsmath}

\begin{document}

\section*{R\'ubrica de Evaluaci\'on - Trabajo de Econometría}

\subsection*{1. Estad\'istica Descriptiva y Base de Datos (15 pts)}
\begin{itemize}
  \item \textbf{Excelente (15-13)}: Estad\'isticas descriptivas claras y completas. Correlaciones relevantes bien interpretadas. Visualizaciones pertinentes y claras. Se documenta el origen y tipo de la base de datos.
  \item \textbf{Aceptable (12-6)}: Estad\'isticas o visualizaciones incompletas o poco discutidas. Correlaciones listadas sin an\'alisis. Menci\'on breve al origen de los datos.
  \item \textbf{Insuficiente (<6)}: An\'alisis descriptivo ausente o confuso. No se indica el origen de los datos ni se visualiza informaci\'on clave.
\end{itemize}

\subsection*{2. Planteamiento del Modelo e Hipótesis de la investigación (20 pts)}

\begin{itemize}

  \item \textbf{Modelo poblacional (8 pts):}
    \begin{itemize}
      \item \textbf{Excelente (6--8 pts):} Fundamenta adecuadamente la elección de la variable dependiente y las explicativas. Plantea el modelo poblacional con justificaciones sólidas desde la teoría económica y apoyado en el análisis descriptivo.
      \item \textbf{Aceptable (3--5 pts):} El modelo está presente pero la justificación es parcial, poco clara o limitada a una sola dimensión (económica o descriptiva).
      \item \textbf{Insuficiente (2--0 pts):} Modelo ausente, incorrecto o con una justificación deficiente o errónea.
    \end{itemize}

  \item \textbf{Planteamiento de hipótesis de investigación (12 pts):}
    \begin{itemize}
      \item \textbf{Excelente (12--9 pts):} Hipótesis formuladas de manera clara, falsables, y directamente relacionadas con el modelo. Muestran coherencia entre teoría, variables y objetivos del estudio.
      \item \textbf{Aceptable (8--5 pts):} Hipótesis presentes pero con formulación incompleta, poco claras o con débil conexión con el modelo.
      \item \textbf{Insuficiente (4--0 pts):} Hipótesis ausentes, mal formuladas (no falseables), genéricas o sin relación con el modelo planteado.
    \end{itemize}

\end{itemize}


\subsection*{3. Verificaci\'on de Supuestos del teorema de Gauss-Markov (25 pts)}
\textbf{Cada subcriterio vale 5 pts}
\textbf{Cada subcriterio vale 5 pts}

\begin{itemize}

  \item \textbf{Exogeneidad:} \\
  \textbf{(5--4 pts)} Analiza adecuadamente la exogeneidad de las variables explicativas, con justificación teórica y/o empírica clara. \\
  \textbf{(3--2 pts)} Mención breve, incompleta o poco argumentada sobre la exogeneidad. \\
  \textbf{(1--0 pts)} No se analiza la exogeneidad o se presentan errores conceptuales.

  \item \textbf{Homoscedasticidad:} \\
  \textbf{(5--4 pts)} Aplica pruebas adecuadas (e.g., Breusch-Pagan, White) y discute correctamente los resultados. \\
  \textbf{(3--2 pts)} Aplica la prueba, pero sin análisis o con interpretación superficial. \\
  \textbf{(1--0 pts)} No aplica prueba alguna o la interpreta incorrectamente.

  \item \textbf{Normalidad:} \\
  \textbf{(5--1 pts)} Evalúa la normalidad de los errores o la normalidad asintótica (según corresponda), utilizando pruebas o herramientas adecuadas (e.g., Shapiro-Wilk, Jarque-Bera, teorema del límite central), con o sin interpretación. \\
  \textbf{(0 pts)} No se evalúa ningún tipo de normalidad ni se justifica su omisión.

  \item \textbf{Multicolinealidad:} \\
  \textbf{(5--4 pts)} Calcula e interpreta correctamente los VIF. \\
  \textbf{(3--2 pts)} Calcula los VIF pero no ofrece interpretación o esta es vaga. \\
  \textbf{(1--0 pts)} No evalúa la multicolinealidad o hay errores metodológicos.

  \item \textbf{Correcciones si aplica:} \\
  \textbf{(5--4 pts)} Aplica correctamente correcciones necesarias (e.g., robustez, transformaciones, reestimación), presentando resultados finales coherentes. \\
  \textbf{(3--2 pts)} Aplica correcciones parciales o no suficientemente justificadas. \\
  \textbf{(1--0 pts)} No aplica correcciones, aun cuando son necesarias.

\end{itemize}


\subsection*{4. Testeo de hip\'otesis de investigaci\'on (15 pts)}
    \begin{itemize}
      \item \textbf{Excelente (15--10 pts)}: Hip\'otesis contrastadas con los resultados del {\bf modelo corregido.} Resultados claramente interpretados.
      \item \textbf{Aceptable (9--5 pts)}: Pruebas realizadas pero con interpretaci\'on superficial o sin v\'inculo expl\'icito con las hip\'otesis planteadas.
      \item \textbf{Insuficiente (4--0 pts)}: Hip\'otesis no testeadas, mal interpretadas o pruebas incorrectas.
    \end{itemize}


\subsection*{5. Conclusiones y Recomendaciones (10 pts)}

\begin{itemize}
  \item \textbf{Excelente (10--8 pts):} Hallazgos presentados de manera clara y coherente, las hipótesis han sido evaluadas explícitamente, se reconocen las limitaciones del análisis y se proponen recomendaciones o mejoras pertinentes.
  \item \textbf{Aceptable (7--5 pts):} Conclusiones presentes pero formuladas de manera vaga, general, o sin conexión clara con los resultados obtenidos.
  \item \textbf{Insuficiente (4--0 pts):} No se presentan conclusiones claras ni recomendaciones; se omiten aspectos clave del análisis.
\end{itemize}

\subsection*{6. Calidad del documento: redacción, ortografía y referencias (15 pts)}

\begin{itemize}
  \item \textbf{Excelente (15--12 pts):} El documento está bien redactado, con redacción clara, sin errores ortográficos, con coherencia argumentativa y formato uniforme. Las referencias están correctamente citadas y completas.
  \item \textbf{Aceptable (11--7 pts):} Presenta algunos errores de redacción u ortografía. Coherencia parcial en la argumentación o pequeñas inconsistencias en el formato de las referencias.
  \item \textbf{Insuficiente (6--0 pts):} Redacción descuidada, múltiples errores ortográficos, falta de coherencia en la exposición y referencias mal citadas o ausentes.
\end{itemize}


\subsection*{6. Presentación Oral (15 pts)}
\textbf{Cada subcriterio vale 5 pts}

\begin{itemize}

  \item \textbf{Claridad y organización:} \\
  \textbf{Excelente (5 pts):} Presentación fluida, lógica y bien estructurada. Se comunican las ideas con claridad y coherencia. \\
  \textbf{Aceptable (4--3 pts):} Presentación con cierta desorganización o falta de claridad en algunos momentos. \\
  \textbf{Insuficiente (2--0 pts):} Presentación confusa, desordenada o difícil de seguir.

  \item \textbf{Uso del lenguaje técnico:} \\
  \textbf{Excelente (5 pts):} Uso adecuado y preciso del vocabulario técnico del área. \\
  \textbf{Aceptable (4--3 pts):} Uso parcial o con errores menores en términos técnicos. \\
  \textbf{Insuficiente (2--0 pts):} Uso incorrecto, confuso o inapropiado del lenguaje técnico.

  \item \textbf{Defensa y respuestas a preguntas:} \\
  \textbf{Excelente (5 pts):} Todo el grupo responde con claridad, seguridad y fundamentos a las preguntas. \\
  \textbf{Aceptable (4--3 pts):} Responden a las preguntas con dudas o falta de profundidad. \\
  \textbf{Insuficiente (2--0 pts):} No responden adecuadamente. 

\end{itemize}


\subsection*{Total: 100 puntos}

\end{document}
